\documentclass[10pt,a4paper]{amsart}
\usepackage[margin=19mm]{geometry}
\usepackage{amsmath,amssymb,mathtools}
\usepackage[scr=rsfs,cal=euler]{mathalfa}
\usepackage{xfrac}
\usepackage[unicode,hidelinks,bookmarks=false]{hyperref}
\newcommand{\ie}{i.e.\ }
\newcommand{\cf}{cf.\ }
\newcommand{\resp}{respectively\ }
\newcommand{\Subsection}{\subsection}
\providecommand{\nobreakdash}{\nobreak\hbox{-}}
\setlength{\emergencystretch}{2em}
\multlinegap0pt
\usepackage{mathtools}
\usepackage{amssymb}
\usepackage{bm}
\newtheorem{proposition}{Proposition}
\DeclareMathOperator{\DV}{DV}
\newcommand*{\MF}{\mathrm{MF}}
\DeclareMathOperator{\OC}{OC}
\renewcommand*{\phi}{\varphi}
\title{Independence questions in a finite\\axiom-schematization of first-order logic}
\author{Benoit Jubin}
\date{Source: arXiv:2202.10383v3 (CC BY 4.0)}
\begin{document}
\maketitle

\noindent\emph{Source and licence.} Independence questions in a finite\\axiom-schematization of first-order logic. Version arXiv:2202.10383v3 (CC BY 4.0), distributed by arXiv under the Creative Commons Attribution 4.0 International licence, http://creativecommons.org/licenses/by/4.0/, as recorded in the arXiv licence metadata for this exact version and shown on the abstract page https://arxiv.org/abs/2202.10383v3. Full licence notice: \texttt{licenses/arxiv-2202.10383-CC-BY-4.0.txt}.\par\smallskip

\noindent\textit{Editorial scope.} Counted unit: the article's proposition prop:ALLcomm (source0.tex lines 1764--1768). The article's complete original proof of it is delivered with it (source0.tex lines 1770--1801, one \texttt{proof} environment of 32 lines). No mathematics is written, repaired or re-derived: the statement and the proof are the article's own lines, in the article's own notation, and no retained line is substituted or re-flowed.  Retained uncounted as the source-local context the proof needs: lines 1695--1697.  Nothing was omitted from any retained span, so the package carries no omission record.  No retained line contains a citation, so the wrapper carries no bibliography and no bibliography entry.  No retained line contains a figure environment, an image-inclusion command or drawing code, so no image is delivered and the package declares no asset dependency.  Editorial formatting only, in addition: an AMS \texttt{amsart} preamble that restates the article's own declarations and macros used by the retained lines, namely the environments \texttt{proposition} on separate counters, together with the standard packages those lines need.  The retained environments print in this document as Proposition 1 in order of appearance, whereas the article prints the same items with the numbering of its own full document.  The complete native source of the article as distributed in the arXiv e-print for this exact version is bundled unchanged as \texttt{sources/122360/source0.tex}.
\begin{proposition}\label{prop:prov}
A scheme provable from a set of supertrue schemes is supertrue.
\end{proposition}

\begin{proposition}\label{prop:ALLcomm}
In the axiom system $\texttt{TMM} \setminus \{\texttt{spec}, \texttt{ALLeq} \}$, the
scheme~\texttt{ALLcomm} cannot be weakened by adding the DV~condition $\DV(x, y)$ or $\DV(y, \phi)$;
in particular, it is independent.
\end{proposition}

\begin{proof}
An instance of \texttt{ALLcomm} is of the form
$\Psi \coloneqq \forall x \forall y \Phi \to \forall y \forall x \Phi$ (with possibly DV conditions)
where $\Phi \in \MF$.
Its $(i, j)$-transforms other than the $(x, y)$-transform and the $(y, x)$-transform affect
both occurrences of $\Phi$ similarly, so they are instances of \texttt{ALLcomm}, hence are true.
Its $(y, x)$-transform is
$\forall x \forall y \Phi^{x \leftarrow y} \to \forall y \forall y \Phi^{x \leftarrow y}$,
which is essentially an instance of \texttt{spec}, hence is true.
Its $(x, y)$-transform is
$\forall x \forall x \Phi^{y \leftarrow x} \to \forall y \forall x \Phi^{y \leftarrow x}$.

An instance of $(\texttt{ALLcomm}, \DV(x, y))$ is of the form $(\Psi, \DV(x, y))$ with~$\Psi$
as above.
Since the $(x, y)$-transform is illegitimate on it, $(\texttt{ALLcomm}, \DV(x, y))$ is supertrue.
An instance of $(\texttt{ALLcomm}, \DV(y, \phi))$ is of the form
$(\Psi, \{ \{y, m \} \mid m \in \OC(\Phi) \})$ with~$\Psi$ as above.
Its $(x, y)$-transform (legitimate or not) is essentially generalization over a nonfree variable,
hence is true.
Therefore, $(\texttt{ALLcomm}, \DV(y, \phi))$ is supertrue.

On the other hand, the instance
$(\forall x \forall y z \equiv t \to \forall y \forall x z \equiv t, \{\{x, z\}, \{x, t\}\})$
of $(\texttt{ALLcomm}, \DV(x, \phi))$ is not supertrue since its $(x, y)$-transform is
$(\forall x \forall x z \equiv t \to \forall y \forall x z \equiv t, \{\{x, z\}, \{x, t\}\})$,
which is not true, as its instance $z \leftarrow y$ shows.

By the above paragraphs, the three previous lemmas, and Proposition~\ref{prop:prov}, all schemes
provable from $\texttt{TMM} \setminus \{\texttt{spec}, \texttt{ALLeq}, \texttt{ALLcomm} \}$ and
the above two weakenings of \texttt{ALLcomm} are supertrue.
Since \texttt{ALLcomm} is not, it is not provable from them.
\end{proof}

\end{document}
